\section{Conclusão}
\label{sec:ia-conclusão}

A inteligência artificial transmutou-se da especulação para uma utilidade pragmática inflexível. Neste compêndio, delineamos de forma aprofundada como os fundamentos matemáticos e as ontologias computacionais conferem alma, tempo e prognóstico ao até então estático conceito de Gêmeos Digitais.

O vetor inegável e já consolidado está na micro-segmentação (e.g., redes nnU-Net), que subverteu fluxos de trabalho arcaicos, inserindo cadência e consistência sobre-humana ao ato morfológico de separar tecido sadio de doente \cite{duggal2026, elsaid2025digital}. As simulações biomédicas multivariáveis — englobando ancoragem ortodôntica, expansões radiculares guiadas \cite{li2024} e otimizações implantares precisas \cite{alshadidi2024} — evidenciam resultados superlativos frente à imaginação analógica do profissional, ainda que clamem por consolidação e validação de longo alcance temporal.

Nada disso exime a classe acadêmica e corporativa do ônus moral e social das disparidades operacionais: explicabilidade XAI como mandamento ético, a auditoria rigorosa de desvios algorítmicos em minorias, e a inadiável integração técnica das soluções baseadas em LLMs e gêmeos com a prática popular unificada (\textit{open-source}) \cite{robles2023}. 

A orquestração através de plataformas abertas, sustentada por pipelines de avaliação cruzada rigorosa sob princípios de Test-Driven Development (TDD), assegurará não apenas a escalabilidade sem precedentes dessas arquiteturas em um mercado bilionário, como essencialmente resguardará o princípio básico da odontologia em seu imperativo de garantir a dignidade humana, transformando a arte de reabilitar os tecidos orais em uma ciência ontológica exata e indiscutivelmente acessível \cite{infante2024, frota2025}.
